\scanpage{14}

\deleted{\uline{Corollaire 2} Soient $C,C'$ deux courbes sur $S$, munies d'isomorphismes de \deleted{schémas} schémas abéliens polarisés $\varphi:\deleted{\unsure{\varphi}}J^0_{C/S}\to A$ et $\varphi':J^0_{C'/S}\to A^{\unsure{*|\prime}}$. \deleted{Alors} On suppose que $S$ est réduit}

\uline{Corollaire 2} (i) Supposons que $S$ soit \inserted{\uline{réduit}, et} \uline{\unsure{connexe|normal}} \unsure{aux|avec} \deleted{$\unsure{s\in S}$} pts $s\in S$ tels que $A_{\deleted{\unsure{s}}s}$ soit une jacobienne \deleted{\unsure{non}} \uline{hyperelliptique}. Pour que $A$ soit un schéma jacobien \inserted{\uline{polarisé}}, il f. et s. que ses fibres le soient ; ou même que les fibres gén. (correspondant aux pts gén. des composantes de $S$) le soient, et que les diviseurs de polarisation des \inserted{\uline{\unsure{ttes|de}} les} fibres puissent être pris irréductibles
% slanted
\margin{il suffit que \deleted{\unsure{l'on}} $S$ soit réduit aux dits pts et que le morphisme $\unsure{\widetilde S|S'|s}\to S$ (\unsure{cf le n° suivant}) soit \uline{non ramifié} en $s$.}
% slanted
\margin{faux sous cette forme}

(ii) Supposons que $S$ soit \uline{réduit} aux \deleted{\unsure{pts}} $s\in S$ tels que $A_s$ soit hyperelliptique. Soient $C,C'$ deux courbes sur $S$, munies d'isom. de schémas abéliens polarisés $\varphi:J^0_{C/S}\mathbin{\unsure{-|\to}}A$, $\varphi':J^0_{C'/S}-A$. Alors il existe un $S$-isomorphisme $u:C-C'$ et un seul \deleted{tel que \unsure{qu}} compatible avec $\varphi$ et $\varphi'$.

% archive sheet number: 13
